\section{Transfer through Witt coefficients}
\label{sec:witt-transfer}

Witt cohomology supplies an injective base-change map with which to
reflect the relation of \cref{thm:geometric-transport} from $Y_C$ to
$Y$. We first prove this injectivity by reducing to trivial Witt
coefficients for a compact group. We then carry the relation to
deformation coefficients and prove nonvanishing there using the
rank-two theory and the retraction onto its Witt constants.

Recall the notation
\[
 G=P_3\times P_2,\qquad
 X=[\Spd\widehat L/P_3^1],\qquad
 H=P_2\times\overline Z,
 \quad \overline Z=P_3/P_3^1.
\]
Thus $Y=X/H$, with the residual action supplied by quotient
descent.  For any of these stacks $Z$, write
\[
 R\Gamma_v(Z,W(\cO^\flat))
   :=\Rlim_m R\Gamma_v(Z,W_m(\cO^\flat)),
 \qquad
 H_W^i(Z):=H^i\bigl((R\Gamma_v(Z,W(\cO^\flat)))(*)\bigr).
\]
The derived global sections are solid complexes; \((*)\) is the exact,
limit-preserving point evaluation, so \(H_W^i(Z)\) is an ordinary abelian
group. This is an integral definition. Inverting $p$ below always takes
place after this derived limit and after cohomology.

\subsection{Witt constants and base change}

In \cite[Remark~3.6.12(2)]{BMR2026}, Barthel and coauthors propose
studying mixed-characteristic Witt cohomology on the unquotiented
punctured Banach--Colmez space. Here the finite norm-unit quotient
$P_3^1/\ker N_3\simeq\mu_{p-1}$ has already removed the higher
characteristic-$p$ contribution in \cref{prop:constant-acyclicity}.
That constants equivalence lifts through the finite Witt levels,
giving the following more restricted comparison.

\begin{proposition}[Witt comparison]\label{prop:witt-comparison}
The natural constant maps give the commutative square of integral
cohomology groups
\begin{equation}\label{eq:witt-square}
\begin{tikzcd}[column sep=large,row sep=large]
 H^i(H,W(k)) \arrow[r,"\sim"] \arrow[d]
   & H_W^i(Y) \arrow[d]
   \\
 H^i(H,W(C^\flat)) \arrow[r,"\sim"]
   & H_W^i(Y_C).
\end{tikzcd}
\end{equation}
The actions of $H$ on the two left-hand coefficient rings are trivial.
The right vertical map is injective after inverting $p$, for every $i$.
Moreover, quotient descent gives a natural identification
\begin{equation}\label{eq:witt-group-identification}
 R\Gamma_v(Y,W(\cO^\flat))
   \simeq R\Gamma(G,W(\widehat L)).
\end{equation}
\end{proposition}

\begin{proof}
For a characteristic-$p$ ring $R$, restriction of Witt components and
Verschiebung give an exact sequence of additive groups
\begin{equation}\label{eq:witt-devissage}
 0\longrightarrow R\xrightarrow{\ V^{m-1}\ }W_m(R)
   \longrightarrow W_{m-1}(R)\longrightarrow0
 \qquad(m\geq2).
\end{equation}
The sequence is natural for ring maps and automorphisms.  We use it as
a sequence of additive coefficient objects; no assertion of linearity
over $R$ is needed.  It also gives an exact sequence of the associated
solid modules and of the corresponding sheaves on the v-site.

\Cref{prop:constant-acyclicity} identifies the actual constant maps
\[
 k\longrightarrow R\Gamma(P_3^1,\widehat L),\qquad
 C^\flat\longrightarrow R\Gamma_v(X_C,\cO^\flat)
\]
as equivalences with their residual $H$-actions.  As proved there,
the characteristic-$p$ vanishing uses only the central scalar
restrictions of \cite[Propositions~3.5.6 and~3.6.9]{BMR2026}, identified
with $T_3$ by the determinant calculation in \cref{prop:tower}, and finite
norm-unit
invariants as in \cite[Lemma~5.2.5]{BMR2026}.  Full $H$-equivariance
comes from the actual maps of constants and quotient descent.
Apply \eqref{eq:witt-devissage} to the source and target of each map.
Induction on $m$, using the resulting morphisms of exact triangles,
gives equivalences
\begin{align*}
 W_m(k)&\xrightarrow{\sim}
       R\Gamma(P_3^1,W_m(\widehat L)),\\
 W_m(C^\flat)&\xrightarrow{\sim}
       R\Gamma_v(X_C,W_m(\cO^\flat)).
\end{align*}
Every map used in this induction is the map induced by constants.
In particular the equivalences are $H$-equivariant, including the
$\overline Z$-action; $H$ fixes the constant rings.

Derived global sections and derived invariants preserve limits.
Passing to $\Rlim_m$ therefore gives the same equivalences with full
Witt coefficients.  Here the ordinary Witt module is also the derived
limit of its truncations: in Witt coordinates, restriction has the
continuous set-theoretic section obtained by appending a zero.
Consequently restriction is surjective on continuous functions from
every profinite parameter space.  The truncation towers have no higher
derived-limit term in the solid coefficient category.  The argument
applies to $k$, $C^\flat$, and $\widehat L$ with their respective
topologies.  Taking residual $H$-invariants gives the horizontal
isomorphisms of \eqref{eq:witt-square}.  Naturality of the constant
maps under $k\to C^\flat$ gives its commutativity.

For \eqref{eq:witt-group-identification}, recall
$\widetilde Y=\Spd\widehat L$ from \cref{sec:completed-tower}.
This affinoid perfectoid space is acyclic for $\cO^\flat$, also with
profinite parameters.  Applying \eqref{eq:witt-devissage} gives, for
every profinite set $S$, the finite equivalence induced by sections
\[
 \alpha_{S,m}:C_{\cts}(S,W_m(\widehat L))[0]
       \xrightarrow{\ \sim\ }
 R\Gamma_v(\widetilde Y\times\underline S,W_m(\cO^\flat)).
\]
Here Witt coordinates identify $W_m(C_{\cts}(S,\widehat L))$ with
$C_{\cts}(S,W_m(\widehat L))$.  The equivalence is natural in $S$
and in the $G$-action.  Write $u^a_{\widetilde Y,G,m}$ for the
degree-$a$ map of \eqref{eq:finite-constant-cochains}.  The actual
finite coefficient units give the commutative square
\begin{equation}\label{eq:finite-witt-constant-square}
\begin{tikzcd}[column sep=large]
 C_{\cts}(G^a,\Lambda_m)[0]
    \arrow[r,"u^a_{\widetilde Y,G,m}"] \arrow[d,"\iota_m"']
   &R\Gamma_v(\widetilde Y\times\underline{G^a},\Lambda_m)
       \arrow[d,"\omega_m"]\\
 C_{\cts}(G^a,W_m(\widehat L))[0]
    \arrow[r,"\alpha_{G^a,m}","\sim"']
   &R\Gamma_v(\widetilde Y\times\underline{G^a},W_m(\cO^\flat)).
\end{tikzcd}
\end{equation}
The left map applies $\Lambda_m=W_m(\F_p)\to W_m(\widehat L)$
pointwise; the right map is induced by the corresponding sheaf unit.
Both composites are the same locally constant Witt section.  These
squares commute with the bar faces, including the action face, and
with restriction in $m$.  The same calculation with an extra
profinite parameter retains the solid structure.  The condensed
standard resolution, quotient descent, and then $\Rlim_m$ therefore
give \eqref{eq:witt-group-identification}.  The full Witt module is
the derived limit of its truncations as shown above.  Under point
evaluation of this identification, the Witt image of
$u_{\widetilde Y,G,\mathrm{int}}$ is exactly the coefficient map
induced by $\Z_p\to W(\widehat L)$ and
\eqref{eq:group-cochain-bridge}.

It remains to prove injectivity, since an equivalence of the horizontal
groups alone would not imply it.  The finite-resolution statement of
\cref{lem:solid-resolutions} applies to $H$.  Indeed an element of
order $p$ in $P_2$ would have minimal polynomial of degree $p-1>2$
over $\Q_p$, whereas the reduced characteristic polynomial in $D_2$
has degree two.  Thus $P_2$, and hence $H=P_2\times\Z_p$, has no
$p$-torsion.
Choose a bounded resolution $P_\bullet\to\Z_p$ by finitely generated
projective $\Z_p[[H]]$-modules.  For a trivial solid coefficient module
$M$, the point-valued cochain complex is
\[
 \Hom_{\Z_p[[H]]}(P_\bullet,M)
   =\Hom_{\Z_p}(Q_\bullet,M),
 \qquad
 Q_\bullet=\Z_p\otimes_{\Z_p[[H]]}P_\bullet.
\]
Every $Q_j$ is finite free over $\Z_p$.  Thus this complex is obtained
from a fixed bounded finite free $\Z_p$-complex by scalar extension
to the underlying module of $M$; its differentials are the same
$\Z_p$-linear matrices.  This description is natural in $M$ and is
valid for the valued Witt topology as well as for the $p$-adic one,
by \cref{lem:solid-resolutions}.

Take $A=W(k)$ and $B=W(C^\flat)$.  The map $A\to B$ is injective.
Since $A$ is a discrete valuation ring with uniformizer $p$ and $B$
is $p$-torsion-free, $B$ is flat over $A$.  Applying the finite
complex just described gives
\begin{equation}\label{eq:witt-scalar-extension}
 H^i(H,A)[1/p]\otimes_{A[1/p]}B[1/p]
   \xrightarrow{\sim}H^i(H,B)[1/p].
\end{equation}
The field extension $A[1/p]\hookrightarrow B[1/p]$ is faithfully
flat, so the natural map from the first cohomology group to its
scalar extension is injective.  The square now proves the assertion.
All infinite limits preceded this finite-complex argument.
\end{proof}

\subsection{Reflecting the geometric relation}

Via \eqref{eq:group-cochain-bridge}, let $c_n^W$ denote the image of
$e_{P_n}$ under projection from $G$ and the coefficient map
$\Z_p\to W(\widehat L)$:
\[
 c_n^W\in H^3(G,W(\widehat L))[1/p]\qquad(n=2,3).
\]
Write $c_n$ for its further image under the natural equivariant map
$W(\widehat L)\to\pi_0\widehat B^{\square}$ of
\cref{prop:completed-theory}.

\begin{proposition}[Transport to deformation coefficients]
\label{prop:coefficient-transport}
For the nonzero scalar $a\in\Q_p^\times$ of
\cref{thm:geometric-transport}, one has
\begin{equation}\label{eq:witt-relation}
 c_3^W=a\,c_2^W
       \quad\hbox{in }H^3(G,W(\widehat L))[1/p],
\end{equation}
and hence
\begin{equation}\label{eq:deformation-relation}
 c_3=a\,c_2
       \quad\hbox{in }H^3(G,\pi_0\widehat B^{\square})[1/p].
\end{equation}
Under \eqref{eq:coefficient-descent}, the preimage of $c_2$ in
\[
 H^3(H,\pi_0E_2^{\square})[1/p]
\]
is obtained by inflating $e_{P_2}$ and then including constants.
\end{proposition}

\begin{proof}
Let $\omega_{Z,m}$ be the map on cohomology complexes induced by
$\Lambda_m=W_m(\F_p)\to W_m(\cO^\flat)$ for $Z=Y$ or $Y_C$.
These maps are compatible with restriction in $m$.  Their derived
inverse limit gives a map
\[
 H^3_{\mathrm b}(Z)\longrightarrow H_W^3(Z)[1/p].
\]
Let $j:Y_C\to Y$ denote base change, and put
\[
 \beta_m=\omega_{Y,m}u_{\widetilde Y,G,m},\qquad
 \beta_{C,m}=\omega_{Y_C,m}u_{\widetilde Y_C,G,m}.
\]
The finite constant maps and the finite Witt units give
\begin{equation}\label{eq:finite-witt-base-change}
\begin{tikzcd}[column sep=large]
 C_{\cts}^{\bullet}(G,\Lambda_m)
      \arrow[r,"\beta_m"] \arrow[d,equal]
   &R\Gamma_v(Y,W_m(\cO^\flat))\arrow[d,"j^*"]\\
 C_{\cts}^{\bullet}(G,\Lambda_m)
      \arrow[r,"\beta_{C,m}"']
   &R\Gamma_v(Y_C,W_m(\cO^\flat)).
\end{tikzcd}
\end{equation}
This square commutes in each cosimplicial degree before totalization,
and is compatible with $q_n:G\to P_n$ and reduction in $m$.
Take the derived coefficient limit, cohomology, and then invert $p$.
The classifying-pullback convention and
\eqref{eq:constant-cochain-naturality} give
$j^*(e_{P_n}|_Y)=e_{P_n}|_{Y_C}$. Moreover,
\eqref{eq:finite-witt-constant-square} identifies the Witt image of
$e_{P_n}|_Y$ with $c_n^W$ under
\eqref{eq:witt-group-identification}. Thus the geometric relation says
that $c_3^W-a c_2^W$ maps to zero in $H_W^3(Y_C)[1/p]$.
The right vertical map of \eqref{eq:witt-square} is injective after
inverting $p$, so this difference is already zero on $Y$. This proves
\eqref{eq:witt-relation}. Applying the coefficient map
$W(\widehat L)\to\pi_0\widehat B^{\square}$ then proves
\eqref{eq:deformation-relation}.

Finally, \eqref{eq:coefficient-descent} is induced by the actual map
$E_2^{\square}\to\widehat B^{\square}$, by first taking
$P_3^1$-invariants and then $H$-invariants.  Its compatibility with
constants is part of \cref{prop:completed-theory}, whose finite Koszul
argument proves the required coefficient instance of the descent
construction in \cite[Proposition~4.6.12]{BMR2026}.
The action of $\overline Z$ on the source is trivial.  Thus
inflation from $P_2$ followed by this map is exactly the coefficient
map defining $c_2$.  This proves the last assertion while retaining
the residual action.
\end{proof}

\subsection{Nonvanishing in the deformation ring}

The reflected relation reduces nonvanishing to $c_2$. Coefficient
descent identifies this class with constants in the rank-two theory;
the additive retraction onto Witt constants detects it there.

\begin{lemma}[Nonvanishing of constants]\label{lem:constant-nonvanishing}
The class $c_2$ is nonzero.  Consequently $c_3$ is nonzero as well.
\end{lemma}

\begin{proof}
Set $A_k=\pi_0E_2^{\square}(*)$.  The class corresponding to $c_2$
under \eqref{eq:coefficient-descent} restricts along $P_2\hookrightarrow H$
to the coefficient image of $e_{P_2}$ in $H^3(P_2,A_k)[1/p]$.
It suffices to prove that this image is nonzero.

Choose an algebraic closure $\overline k$ of $k$ and let
$A_{\overline k}$ be the deformation ring of the same height-two
formal group after this base extension.  Deformation theory gives
a natural continuous $P_2$-equivariant map
\[
 A_k\longrightarrow A_{\overline k}.
\]
In compatible coordinates these rings are $W(k)[[u'_1]]$ and
$W(\overline k)[[u'_1]]$, with their $(p,u'_1)$-adic topologies;
equivariance follows from deformation base change, without requiring
the coordinate to be fixed.  Since $P_2$ is nonextended, it fixes
$W(\overline k)$ pointwise.

The splitting theorem of Barthel--Schlank--Stapleton--Weinstein
\cite[Proposition~2.5.1 and Corollary~2.5.9]{BSSW2025} gives a
continuous additive equivariant retraction
\[
 r:A_{\overline k}\longrightarrow W(\overline k)
\]
of the inclusion of Witt constants. Restrict its equivariance to $P_2$.
Additivity and continuity make it $\Z_p$-linear. After condensation
it is a morphism of solid $P_2$-modules, and hence
induces a retraction on the continuous derived invariants used here.

The composite of the coefficient maps
\[
 \Z_p\longrightarrow A_k\longrightarrow A_{\overline k}
       \xrightarrow{r}W(\overline k)
\]
is the usual inclusion of constants.  Applying the finite free
complex argument from \cref{prop:witt-comparison} with $P_2$ in
place of $H$ gives
\[
 H^3(P_2,W(\overline k))[1/p]
   \cong H^3_{\mathrm b}(P_2)
          \otimes_{\Q_p}W(\overline k)[1/p].
\]
The image of $e_{P_2}$ here is $e_{P_2}\otimes1$, which is nonzero
by \cref{prop:group-classes}.  If its image in
$H^3(P_2,A_k)[1/p]$ vanished, the displayed composite would also
send it to zero, a contradiction.  Restriction to $P_2$ therefore
proves $c_2\ne0$, and \eqref{eq:deformation-relation} with $a\ne0$
proves $c_3\ne0$.
\end{proof}

We have thus identified the image of the degree-three group class
under the completed coefficient map.  The remaining step is to
detect it in the homotopy of an ordinary $K(2)$-local spectrum.
